\section*{अध्याय \ref{ch:g7:priorities} --- संक्रियाओं की प्राथमिकता}

\begin{solution}{exo:g7:priorities:1}
$5 + 3 \times 6 = 5 + 18 = 23$।

$(5 + 3) \times 6 = 8 \times 6 = 48$।

$30 - 12 \div 4 = 30 - 3 = 27$।

$(30 - 12) \div 4 = 18 \div 4 = 4.5$।
\end{solution}

\begin{solution}{exo:g7:priorities:2}
$24 - 6 + 2 = 18 + 2 = 20$ (बाएँ से दाएँ)।

$24 \div 6 \times 2 = 4 \times 2 = 8$ (बाएँ से दाएँ)।

$24 - 6 \times 2 = 24 - 12 = 12$ (पहले गुणा)।

$24 \div (6 \times 2) = 24 \div 12 = 2$।
\end{solution}

\begin{solution}{exo:g7:priorities:3}
$7 + 3 \times (10 - 6) = 7 + 3 \times 4 = 7 + 12 = 19$।

$40 - (2 + 3) \times 4 = 40 - 5 \times 4 = 40 - 20 = 20$।

$60 \div \bigl(2 \times (7 - 4)\bigr) = 60 \div (2 \times 3)
= 60 \div 6 = 10$।
\end{solution}

\begin{solution}{exo:g7:priorities:4}
$\dfrac{18 + 6}{8} = \dfrac{24}{8} = 3$; \quad
$\dfrac{30}{7 - 2} = \dfrac{30}{5} = 6$; \quad
$\dfrac{5 \times 8}{4 + 6} = \dfrac{40}{10} = 4$।
\end{solution}

\begin{solution}{exo:g7:priorities:5}
$8 \times 102 = 8 \times (100 + 2) = 800 + 16 = 816$।

$25 \times 99 = 25 \times (100 - 1) = 2500 - 25 = 2475$।

$14 \times 12 + 14 \times 8 = 14 \times (12 + 8) = 14 \times 20 = 280$।

$37 \times 45 - 37 \times 35 = 37 \times (45 - 35) = 37 \times 10 =
370$।
\end{solution}

\begin{solution}{exo:g7:priorities:6}
``$50$ का और $6$ तथा $7$ के \omterm{def:g2:mult:def}{गुणनफल} का \omterm{ex:g1:subtraction:difference}{अंतर}'':
$50 - 6 \times 7 = 50 - 42 = 8$ (कोष्ठक की ज़रूरत नहीं, प्राथमिकता
ही काम कर देती है)।

``$4$ और $5$ के \omterm{def:g1:addition:def}{योगफल} का, $10$ और $8$ के \omterm{ex:g1:subtraction:difference}{अंतर} से
\omterm{def:g2:mult:def}{गुणनफल}'': $(4 + 5) \times (10 - 8) = 9 \times 2 = 18$।
\end{solution}

\begin{solution}{exo:g7:priorities:7}
$3 \times 2.50 + 2 \times 1.20 = 7.50 + 2.40 = 9.90$ यूरो।
\end{solution}

\begin{solution}{exo:g7:priorities:8}
(a) कोष्ठक नहीं: $4 + 2 \times 5 - 1 = 4 + 10 - 1 = 13$।

(b) $(4 + 2) \times 5 - 1 = 30 - 1 = 29$।

(c) $(4 + 2) \times (5 - 1) = 6 \times 4 = 24$।
\end{solution}

\begin{solution}{exo:g7:priorities:9}
$k = 1$: $12 + 4 \times 1 = 16$ और $16 \times 1 = 16$ --- ये मिल
जाते हैं (इसीलिए $k = 1$ बुरी जाँच है!)। $k = 2$:
$12 + 4 \times 2 = 20$ पर $16 \times 2 = 32$: अलग। तो दोनों
व्यंजक आम तौर पर बराबर \emph{नहीं} हैं। ग़लती
$12 + 4 \times k = 16 \times k$, $\times$ की प्राथमिकता अनदेखी
करके पहले $12 + 4$ जोड़ने से आती है।
\end{solution}

\begin{solution}{exo:g7:priorities:10}
\emph{1.} समूह का दाम: $12 + 6 \times n$। $n = 5$ के लिए:
$12 + 30 = 42$ यूरो।

\emph{2.} अलग-अलग दाम: $9n$। तुलना: $n = 3$: $30$ बनाम $27$
(अलग सस्ता); $n = 4$: $36$ बनाम $36$ (बराबर); $n = 5$: $42$ बनाम
$45$ (समूह सस्ता)। $5$ लोगों से आगे समूह वाली दर जीतती है
(हर अतिरिक्त व्यक्ति $9$ की जगह $6$ का पड़ता है, इसलिए फ़ासला बढ़ता ही जाता है)।
\end{solution}

\begin{solution}{exo:g7:priorities:11}
संभव उत्तर (कई हैं):
\[
24 = 1 \times 2 \times 3 \times 4, \qquad
10 = 1 + 2 + 3 + 4, \qquad
1 = (4 - 3) \times (2 - 1).
\]
\end{solution}

\begin{solution}{pb:g7:priorities:1}
\textbf{1.} $7 \times 102 = 7 \times (100 + 2) = 700 + 14 =
714$; \quad $9 \times 98 = 9 \times (100 - 2) = 900 - 18 = 882$।

\textbf{2.} $17 \times 13 + 17 \times 7 = 17 \times (13 + 7) =
17 \times 20 = 340$; \quad $45 \times 101 - 45 = 45 \times (101
- 1) = 45 \times 100 = 4\,500$।

\textbf{3.} $5 \times 87 = 870 \div 2 = 435$। जायज़ क्यों:
$5 = 10 \div 2$, इसलिए $5$ से गुणा करना यानी $10$ से गुणा करके
$2$ का भाग देना (एक ही स्तर पर $\times$ और $\div$ का बाएँ से
दाएँ क्रम बाक़ी सँभाल लेता है)।

\textbf{4.} $25 \times 32 = 3\,200 \div 4 = 800$, क्योंकि
$25 = 100 \div 4$। और
$99 \times 47 = (100 - 1) \times 47 = 4\,700 - 47 = 4\,653$:
\omterm{ex:g1:subtraction:difference}{अंतर} पर वितरण नियम
(\cref{thm:g7:priorities:distributivity})।

\textbf{5.} $8 \times (5 + 3) = 64$ और $8 \times 5 + 8 \times 3
= 40 + 24 = 64$: दोनों मिलते हैं --- यही वितरण नियम है। पर
$8 \times (5 \times 3) = 8 \times 15 = 120$, जबकि
$(8 \times 5) \times (8 \times 3) = 40 \times 24 = 960$: नहीं
मिलते। गुणनखंड \emph{\omterm{def:g1:addition:def}{योगफल} के पदों} पर फैलता है, \omterm{def:g2:mult:def}{गुणनफल} के
गुणनखंडों पर नहीं।

\textbf{6.}
\begin{center}
\renewcommand{\arraystretch}{1.15}
\begin{tabular}{ccl}
$13$ & $24$ & रखी ($13$ विषम) \\
$6$ & $48$ & काटी ($6$ सम) \\
$3$ & $96$ & रखी \\
$1$ & $192$ & रखी \\
\end{tabular}
\end{center}
बची हुई दाईं संख्याओं का \omterm{def:g1:addition:def}{योगफल}: $24 + 96 + 192 = 312$, और
सचमुच $13 \times 24 = 312$।

\textbf{7.} $21$ को \omterm{def:g3:division:half}{आधा} करना: $21, 10, 5, 2, 1$; $17$ को दुगुना:
$17, 34, 68, 136, 272$। सम बाएँ $10$ और $2$ कट जाते हैं;
बची हुई $17 + 68 + 272 = 357$। जाँच:
$21 \times 17 = 357$।

\textbf{8.} $16$ को \omterm{def:g3:division:half}{आधा} करना: $16, 8, 4, 2, 1$; $25$ को दुगुना:
$25, 50, 100, 200, 400$। आख़िरी $1$ को छोड़कर हर बाईं संख्या
सम है: केवल एक पंक्ति बचती है, और उत्तर $400$ है। शुद्ध दुगुनों
से बना बायाँ स्तंभ केवल अपनी आख़िरी पंक्ति रखता है --- विधि
सिमटकर केवल दुगुना करना रह जाती है।

\textbf{9.} $13 = 12 + 1$, इसलिए वितरण नियम से
$13 \times 24 = 12 \times 24 + 1 \times 24$; और
$12 \times 24 = 6 \times 48$, क्योंकि एक गुणनखंड को \omterm{def:g3:division:half}{आधा} और दूसरे
को दुगुना करने से \omterm{def:g2:mult:def}{गुणनफल} नहीं बदलता
($12 \times 24 = 6 \times 2 \times 24 = 6 \times 48$)। विषम
बाईं संख्या $13$ ने दाईं संख्या की एक प्रति अलग कर दी: पहली
पंक्ति $24$ जमा करती है।

\textbf{10.} $6 \times 48 = 3 \times 96$ (आधा-दुगुना, कुछ जमा
नहीं: $6$ सम है)। $3 \times 96 = 2 \times 96 + 96 =
1 \times 192 + 96$: विषम $3$, $96$ जमा करता है। आख़िर में
$1 \times 192 = 192$: विषम $1$, $192$ जमा करता है और सारणी
ख़त्म। कुल जमा: $24 + 96 + 192 = 312$। कोई पंक्ति अपनी दाईं
संख्या ठीक तभी जमा करती है जब उसकी बाईं संख्या विषम हो --- यानी
विधि ``विषम बाईं वाली पंक्तियाँ रखो और जोड़ो'' ठीक वितरण नियम का
हिसाब-किताब है।

\textbf{11.} $1 + 4 + 8 = 13$। $21$ के लिए बची हुई पंक्तियाँ
$17$, $68 = 4 \times 17$ और $272 = 16 \times 17$ वाली थीं: यानी
दुगुने $1 + 4 + 16 = 21$।

\textbf{12.} $25 = 16 + 8 + 1$ और $42 = 32 + 8 + 2$।

\textbf{13.} लक्ष्य से बड़ी न होने वाली सबसे बड़ी दुगुनी संख्या
लेकर घटा दो; जो बचता है वह छोटा है, और अभी-अभी इस्तेमाल की गई
दुगुनी संख्या से भी छोटा (नहीं तो उससे बड़ी समा जाती)। दोहराने
पर \omterm{def:g3:division:remainder}{शेषफल} लगातार सिकुड़ता है और $0$ तक पहुँचना ही पड़ता है:
प्रक्रिया रुक जाती है, और इस्तेमाल हुए दुगुने सब अलग-अलग होते
हैं। $100$ के लिए: $64$ समाता है ($36$ बचा), $32$ समाता है
($4$ बचा), $4$ समाता है ($0$ बचा): $100 = 64 + 32 + 4$।

\textbf{14.} $35$ की दुगुनी सारणी: $1 \to 35$, $2 \to 70$,
$4 \to 140$, $8 \to 280$, $16 \to 560$। चूँकि
$19 = 16 + 2 + 1$, वही पंक्तियाँ चुनो:
$560 + 70 + 35 = 665$। जाँच: $19 \times 35 = 665$।

\textbf{15.} वही तीन जो लिपिक और किसान के थे:
$24 + 96 + 192$ --- $24$ के वे दुगुने जो
$13 = 1 + 4 + 8$ ने चुने। संख्याओं को दुगुनों के \omterm{def:g1:addition:def}{योगफल} के रूप
में लिखना ठीक उन्हें $0$ और $1$ से लिखना है; पूरी कहानी
\cref{pb:g8:powers:1} में है।

\end{solution}
